\documentclass[12pt, a4paper]{article}

% --- 常用宏包 ---
\usepackage[UTF8]{ctex}
\usepackage{amsmath, amsthm, amssymb, amsfonts}
\usepackage{geometry}
\usepackage{fancyhdr}
\usepackage{enumerate}
\usepackage{graphicx}
\usepackage{hyperref}
\usepackage{bm}
\usepackage{tcolorbox}

% --- 页面设置 ---
\geometry{left=2.5cm, right=2.5cm, top=2.5cm, bottom=2.5cm}
\pagestyle{fancy}
\fancyhf{}
\rhead{\today}
\lhead{近世代数作业}
\cfoot{\thepage}

% --- 环境定义 ---
\newtcolorbox{problem}[1]{
    colback=gray!5!white,
    colframe=gray!75!black,
    fonttitle=\bfseries,
    title=题目 #1
}

\newenvironment{solution}{
    \begin{proof}[\textbf{解}]
}{
    \end{proof}
}

% --- 个人信息 ---
\title{近世代数作业 12}
\author{李睿阳 (524120910059)}
\date{\today}

\begin{document}
\maketitle


\begin{problem}{1}
    设有域扩张 $F \subset E$，$\alpha, \beta \in E$ 在 $F$ 上代数。设 $p_{\alpha}(x), p_{\beta}(x)$ 分别是 $\alpha, \beta$ 在 $F$ 上的极小多项式，令 $d_{\alpha} = \deg(p_{\alpha}), d_{\beta} = \deg(p_{\beta})$，且 $(d_{\alpha}, d_{\beta}) = 1$。证明：$[F(\alpha, \beta) : F] = d_{\alpha} \cdot d_{\beta}$，从而有 $p_{\alpha}(x)$ 是 $F(\beta)[x]$ 中的不可约多项式。
\end{problem}
\begin{solution}
    因为 $p_{\alpha}(x)$ 是 $\alpha$ 在 $F$ 上的极小多项式，有 $[F(\alpha):F] = \deg(p_{\alpha}) = d_{\alpha}$。同理 $[F(\beta):F] = d_{\beta}$。\\
    由扩域次数的乘法原理，$[F(\alpha, \beta):F] = [F(\alpha, \beta):F(\beta)][F(\beta):F]$。\\
    由于 $\alpha$ 是 $F[x]$ 中 $p_{\alpha}(x)$ 的根，且 $F[x] \subseteq F(\beta)[x]$，有 $\alpha$ 在 $F(\beta)$ 上的极小多项式 $q(x)$ 整除 $p_{\alpha}(x)$。\\
    因此 $[F(\alpha, \beta):F(\beta)] = \deg(q) \leq \deg(p_{\alpha}) = d_{\alpha}$。\\
    从而 $[F(\alpha, \beta):F] = [F(\alpha, \beta):F(\beta)] \cdot d_{\beta} \leq d_{\alpha} d_{\beta}$。\\
    因为 $F(\alpha) \subseteq F(\alpha, \beta)$ 且 $F(\beta) \subseteq F(\alpha, \beta)$，得 $d_{\alpha} \mid [F(\alpha, \beta):F]$ 且 $d_{\beta} \mid [F(\alpha, \beta):F]$。\\
    由于 $(d_{\alpha}, d_{\beta}) = 1$，由整除性质得 $d_{\alpha} d_{\beta} \mid [F(\alpha, \beta):F]$。\\
    综合以上两方面得 $[F(\alpha, \beta):F] = d_{\alpha} d_{\beta}$。\\
    由 $d_{\alpha} d_{\beta} = [F(\alpha, \beta):F(\beta)] \cdot d_{\beta}$ 得到 $[F(\alpha, \beta):F(\beta)] = d_{\alpha}$。\\
    即 $\alpha$ 在 $F(\beta)$ 上的极小多项式次数为 $d_{\alpha}$，又 $p_{\alpha}(x)$ 在 $F(\beta)[x]$ 中且以 $\alpha$ 为根，所以 $p_{\alpha}(x)$ 为 $\alpha$ 在 $F(\beta)$ 上的极小多项式，从而 $p_{\alpha}(x)$ 在 $F(\beta)[x]$ 中不可约。
\end{solution}


\begin{problem}{2}
    求 $(x^2 - 2)(x^3 - 2)$ 在 $\mathbb{Q}$ 上的分裂域。
\end{problem}
\begin{solution}
    设多项式 $f(x) = (x^2 - 2)(x^3 - 2)$。\\
    $x^2 - 2$ 在 $\mathbb{Q}$ 上的根为 $\pm \sqrt{2}$。\\
    $x^3 - 2$ 在 $\mathbb{Q}$ 上的根为 $\sqrt[3]{2}, \sqrt[3]{2}\omega, \sqrt[3]{2}\omega^2$，其中 $\omega = e^{2\pi i/3} = \frac{-1 + \sqrt{3}i}{2}$。\\
    分裂域 $K = \mathbb{Q}(\sqrt{2}, -\sqrt{2}, \sqrt[3]{2}, \sqrt[3]{2}\omega, \sqrt[3]{2}\omega^2) = \mathbb{Q}(\sqrt{2}, \sqrt[3]{2}, \omega)$。\\
    由 $\omega = \frac{-1 + \sqrt{3}i}{2}$，得 $\mathbb{Q}(\omega) = \mathbb{Q}(\sqrt{3}i)$。\\
    因此分裂域为 $K = \mathbb{Q}(\sqrt{2}, \sqrt[3]{2}, \sqrt{3}i)$。
\end{solution}


\begin{problem}{3}
    设 $F$ 是有限域，$\alpha_1, \cdots, \alpha_n \in \bar{F}$。证明：存在一个元素 $\alpha \in F(\alpha_1, \cdots, \alpha_n)$，使得 $F(\alpha_1, \cdots, \alpha_n) = F(\alpha)$。\\
    提示：先证明 $[F(\alpha_1, \cdots, \alpha_n) : F]$ 有限，得到前者是有限域，由有限域的乘法群是循环群，找出 $\alpha$。
\end{problem}
\begin{solution}
    设 $K = F(\alpha_1, \cdots, \alpha_n)$。由于每个 $\alpha_i$ 都在 $\bar{F}$ 中，有 $\alpha_i$ 在 $F$ 上均代数。\\
    有限个代数元生成的扩张是有限扩张，于是 $[K:F]$ 有限。\\
    因为 $F$ 是有限域，且 $K$ 是 $F$ 的有限扩张，从而 $K$ 也是一个有限域。\\
    由有限域的性质，有限域 $K$ 的非零元构成的乘法群 $K^*$ 是一个循环群。\\
    设 $\alpha$ 是 $K^*$ 的一个生成元，则 $K^* = \langle \alpha \rangle$。\\
    即 $K = \{0\} \cup \{\alpha, \alpha^2, \cdots, \alpha^{|K|-1}=1\}$。\\
    有 $F(\alpha) \subseteq K$。又因为 $\alpha \in K$，且 $K$ 的任何非零元素都可以表示为 $\alpha$ 的幂，得到 $K \subseteq F(\alpha)$。\\
    综上所述，$K = F(\alpha)$。
\end{solution}
    

\end{document}
